% GENERATED FILE. Edit canonical lessons, then run npm run generate:books.
% canonical-source-hash: fnv1a64:ed4cbcd621feefaa
% canonical-lessons: GE-C162-sofa, GE-C162-radio, GE-C162-wetter, GE-C162-wind, GE-C162-gewitter

\chapter{A sofa, A radio, The weather, The wind, A thunderstorm}
\label{ch:ge-a-sofa-a-radio-the-weather-the-wind-a-thunderstorm}
\hlchaptermodality{162}

\begin{chapteropening}
\textbf{By the end of this chapter:} \emph{I can say a sofa, a radio, the weather, the wind and a thunderstorm.}

Complete the last lesson of chapter 162: I can say a sofa, a radio, the weather, the wind and a thunderstorm.
\end{chapteropening}

\section[das Sofa]{das Sofa — a sofa}

\label{lesson:GE-C162-sofa}



\begin{quote}
Before the new one: say the German for a balcony, then the German for a cellar.
\end{quote}

\subsection*{You'll want to know: das Sofa}
\textbf{das Sofa} — \textquotedblleft{}a sofa\textquotedblright{}.

\begin{cousinweb}
\textbf{Sofa} came in from French \emph{sofa}, which EtymDB 2.1 traces to Arabic \emph{ṣuffa}, \textquotedblleft{}a long raised bench\textquotedblright{} — borrowed, like English \emph{sofa}.
\end{cousinweb}

\begin{grammarlens}[title={What you've built}]
One more word for this chapter.
\end{grammarlens}

\subsection*{Your turn}
\begin{itemize}
\raggedright
  \item \emph{Say it:} \emph{das Sofa}
  \item \emph{Say it:} \emph{das Sofa}, once more
  \item \emph{Say it:} say \emph{das Sofa}
  \item \emph{Recall:} say the German for a balcony, then the German for a cellar, then say \emph{das Sofa} again
\end{itemize}

\begin{tcolorbox}[breakable,colback=teal!4,colframe=teal!35!black,title={Before you move on}]
What does das Sofa mean? (\textquotedblleft{}A sofa\textquotedblright{}.) Say it once more.
\end{tcolorbox}
\section[das Radio]{das Radio — a radio}

\label{lesson:GE-C162-radio}



\begin{quote}
Before the new one: say the German for a cellar, then the German for a sofa.
\end{quote}

\subsection*{You'll want to know: das Radio}
\textbf{das Radio} — \textquotedblleft{}a radio\textquotedblright{}.

\begin{cousinweb}
\textbf{Radio} $\leftarrow$ English \emph{radio} (EtymDB 2.1), short for \emph{radiotelegraphy} and built on Latin \emph{radius}, \textquotedblleft{}ray, spoke\textquotedblright{} — borrowed.
\end{cousinweb}

\begin{grammarlens}[title={What you've built}]
One more word for this chapter.
\end{grammarlens}

\subsection*{Your turn}
\begin{itemize}
\raggedright
  \item \emph{Say it:} \emph{das Radio}
  \item \emph{Say it:} \emph{das Radio}, once more
  \item \emph{Say it:} say \emph{das Radio}
  \item \emph{Recall:} say the German for a cellar, then the German for a sofa, then say \emph{das Radio} again
\end{itemize}

\begin{tcolorbox}[breakable,colback=teal!4,colframe=teal!35!black,title={Before you move on}]
What does das Radio mean? (\textquotedblleft{}A radio\textquotedblright{}.) Say it once more.
\end{tcolorbox}
\section[das Wetter]{das Wetter — the weather}

\label{lesson:GE-C162-wetter}



\begin{quote}
Before the new one: say the German for a sofa, then the German for a radio.
\end{quote}

\subsection*{You'll want to know: das Wetter}
\textbf{das Wetter} — \textquotedblleft{}the weather\textquotedblright{}.

\begin{cousinweb}
\textbf{Wetter} $\leftarrow$ Old High German \emph{wetar} $\leftarrow$ Proto-Germanic \emph{*wedrą} — inherited, the same word as English \emph{weather}. (EtymDB 2.1)
\end{cousinweb}

\begin{grammarlens}[title={What you've built}]
One more word for this chapter.
\end{grammarlens}

\subsection*{Your turn}
\begin{itemize}
\raggedright
  \item \emph{Say it:} \emph{das Wetter}
  \item \emph{Say it:} \emph{das Wetter}, once more
  \item \emph{Say it:} say \emph{das Wetter}
  \item \emph{Recall:} say the German for a sofa, then the German for a radio, then say \emph{das Wetter} again
\end{itemize}

\begin{tcolorbox}[breakable,colback=teal!4,colframe=teal!35!black,title={Before you move on}]
What does das Wetter mean? (\textquotedblleft{}The weather\textquotedblright{}.) Say it once more.
\end{tcolorbox}
\section[der Wind]{der Wind — the wind}

\label{lesson:GE-C162-wind}



\begin{quote}
Before the new one: say the German for a radio, then the German for the weather.
\end{quote}

\subsection*{You'll want to know: der Wind}
\textbf{der Wind} — \textquotedblleft{}the wind\textquotedblright{}.

\begin{cousinweb}
\textbf{Wind} $\leftarrow$ Old High German \emph{wint} $\leftarrow$ Proto-Germanic \emph{*windaz} — inherited, the same word as English \emph{wind}. (EtymDB 2.1)
\end{cousinweb}

\begin{grammarlens}[title={What you've built}]
One more word for this chapter.
\end{grammarlens}

\subsection*{Your turn}
\begin{itemize}
\raggedright
  \item \emph{Say it:} \emph{der Wind}
  \item \emph{Say it:} \emph{der Wind}, once more
  \item \emph{Say it:} say \emph{der Wind}
  \item \emph{Recall:} say the German for a radio, then the German for the weather, then say \emph{der Wind} again
\end{itemize}

\begin{tcolorbox}[breakable,colback=teal!4,colframe=teal!35!black,title={Before you move on}]
What does der Wind mean? (\textquotedblleft{}The wind\textquotedblright{}.) Say it once more.
\end{tcolorbox}
\section[das Gewitter]{das Gewitter — a thunderstorm}

\label{lesson:GE-C162-gewitter}



\begin{quote}
Before the new one: say the German for the weather, then the German for the wind.
\end{quote}

\subsection*{You'll want to know: das Gewitter}
\textbf{das Gewitter} — \textquotedblleft{}a thunderstorm\textquotedblright{}.

\begin{cousinweb}
\textbf{Gewitter} $\leftarrow$ Old High German \emph{giwitiri}: \emph{ge-}, the gathering prefix, + \emph{Wetter}, \textquotedblleft{}weather\textquotedblright{} (EtymDB 2.1 has \emph{giwitiri}). A thunderstorm is weather all at once.
\end{cousinweb}

\begin{grammarlens}[title={What you've built}]
One more word for this chapter.
\end{grammarlens}

\subsection*{Your turn}
\begin{itemize}
\raggedright
  \item \emph{Say it:} \emph{das Gewitter}
  \item \emph{Say it:} \emph{das Gewitter}, once more
  \item \emph{Say it:} say \emph{das Gewitter}
  \item \emph{Recall:} say the German for the weather, then the German for the wind, then say \emph{das Gewitter} again
\end{itemize}

\begin{tcolorbox}[breakable,colback=teal!4,colframe=teal!35!black,title={Before you move on}]
What does das Gewitter mean? (\textquotedblleft{}A thunderstorm\textquotedblright{}.) Say it once more.
\end{tcolorbox}
